% status: FULL
\chapter{Reading Shapes, Bytes and Performance Bounds}\label{app:notation}
\begin{ChapterIntent}
Equations are easier to read when every symbol has a job. This chapter
turns the notation used so far into a practical method: identify the
object, write its shape, attach units, then calculate. We will use that
method to count cache bytes and estimate a step's minimum time.
\end{ChapterIntent}

\section{Dimensions and rates}
\begin{longtable}{@{}p{.19\linewidth}p{.48\linewidth}p{.24\linewidth}@{}}
\toprule Symbol & Meaning & Unit\\\midrule
$L,d,d_{\rm ff}$ & Layers, model width, feed-forward width & Counts\\
$H_q,H_{kv},d_h$ & Query heads, KV heads, head width & Counts\\
$V$ & Vocabulary size & Vocabulary entries\\
$B,S$ & Concurrent sequences; context length & Sequences; tokens\\
$P,P_{\rm act}$ & Total and active parameters & Parameters\\
$b_w,b_{kv}$ & Effective weight and cache element size & Bytes/element\\
$q_w=8b_w$ & Effective weight precision & Bits/weight\\
$\beta$ & Bandwidth across a named boundary & Bytes/s\\
$\pi$ & Rate for a named precision and operation & FLOP/s\\
$I$ & Intensity at that boundary & FLOP/byte\\
\bottomrule
\end{longtable}
Four-bit codes have half a byte of payload per element, but effective
$b_w$ includes relevant metadata and mixed types. The early construction writes
$V_{\rm vocab}$ for vocabulary size to distinguish it from the value
matrix $\mathbf{V}$. A locally defined paging block size is not batch $B$.

\section{Shape conventions}
With row-vector activations,
$\mathbf{X}\in\mathbb{R}^{M\times d}$ and
$\mathbf{W}\in\mathbb{R}^{d\times m}$ produce
$\mathbf{XW}\in\mathbb{R}^{M\times m}$. Here $M$ counts work rows, not
necessarily requests: a packed batch may contain several prompt positions
from one request. A bias broadcasts over rows. A transpose changes
orientation; strides describe storage.

For one attention head with $T_q$ query and $T_k$ key positions,
\[
 \mathbf{Q}\in\mathbb{R}^{T_q\times d_h},\quad
 \mathbf{K},\mathbf{V}\in\mathbb{R}^{T_k\times d_h},\quad
 \mathbf{QK}^{\mathsf T}\in\mathbb{R}^{T_q\times T_k}.
\]
Value width is chosen equal to head width here; architectures can differ.
Batch and head dimensions can be explicit or packed. State the kernel's
layout instead of assuming all tensors are contiguous.

\section{A complete cache unit check}
Take an illustrative two-layer model, two KV heads per layer, head width
four and two bytes per cached element. Keys and values together cost
\[
 c_{kv}=2LH_{kv}d_hb_{kv}
       =2\cdot2\cdot2\cdot4\cdot2=64\text{ bytes/token}.
\]
For three requests with five cached tokens each, logical payload is
$3\cdot5\cdot64=960$ bytes. It excludes rounding, tables and workspace.
If blocks hold four tokens, each request needs two blocks, reserving
$3\cdot2\cdot4\cdot64=1536$ bytes. The difference, 576 bytes, is internal
fragmentation, not additional model state.

Zero-based position 5 divided by block size 4 gives quotient and remainder
$(1,1)$. A block table $[7,2]$ selects physical block 2 and token offset 1.
Logical order need not follow physical addresses. Chapter
\ref{ch:paged-kv} builds the translation.

\section{Units and bounds}
GB means $10^9$ bytes and GiB means $2^{30}$ bytes. GB/s is not Gbit/s.
For an illustrative step with $D=2\cdot10^9$ bytes,
$F=8\cdot10^9$ FLOPs, $\beta=100\cdot10^9$ bytes/s and
$\pi=10^{12}$ FLOP/s, an ideal overlapped lower bound is
\[
 t_{\rm step}\geq\max(D/\beta,F/\pi)
              =\max(20,8)\text{ ms}=20\text{ ms}.
\]
Intensity is 4 FLOP/byte; ridge intensity is $\pi/\beta=10$ FLOP/byte.
Four useful output tokens per step imply a ceiling of 200 output tokens/s,
not a measured rate. Additional traffic, launch cost and imperfect overlap
make the actual result slower when the counted work is unavoidable.

TTFT starts at a declared boundary and ends at the first delivered token.
Inter-token latency measures gaps between delivered tokens. Mean time per
output token commonly excludes the first-token interval and is undefined
for a one-token response under that convention. Chapter
\ref{app:reproducing} gives explicit timestamps.

\section{Read an equation aloud}
Start with $\mathbf y=\mathbf x\mathbf W$. Read it as: one row of input
features is transformed into one row of output features by a matrix.
Then write the shapes. If $\mathbf x$ has four entries and
$\mathbf W$ has four rows and three columns, $\mathbf y$ has three
entries. The four input entries are the dimension being reduced.

The component form exposes the calculation:
\[
 y_j=\sum_{i=0}^{3}x_iW_{ij},\qquad j\in\{0,1,2\}.
\]
The index $i$ disappears because it is summed over. The index $j$
remains because it identifies each output coordinate. This is a
practical way to catch a transpose mistake before running code.

Now stack five input rows into $\mathbf X\in\mathbb R^{5\times4}$.
The result has shape $5\times3$. We reused the same transformation
five times. There are still twelve matrix parameters, not sixty.
Parameters belong to the model; rows belong to the current work.
Confusing those counts leads to incorrect memory estimates.

\section{A tensor axis needs a meaning}
An array of shape $[2,3,4]$ contains 24 values, but that fact alone
does not tell us what it represents. The axes might mean batch,
sequence and features, or sequence, heads and head width. An operator
can accept the same physical shape and apply a different computation
if the semantic axes are wrong.

Write the axis names beside the shape when designing an interface.
For attention, distinguish query positions from key positions. During
full prompt evaluation they may have the same count. During ordinary
incremental decode there is one new query position but many key
positions. Reusing a single ambiguous sequence-length variable can
hide that difference.

Shape also does not specify layout. Strides tell us how coordinate
changes map to addresses. A transpose can exchange logical axes
without rearranging storage. Chapter~\ref{app:tensors} established
that distinction; the same discipline applies to a GPU kernel's
layout contract.

\section{Count one object before multiplying by the fleet}
The cache equation is easiest to derive one factor at a time.
One layer stores keys and values, giving a factor of two. Each has
$H_{kv}$ heads, each head has $d_h$ elements, and each element uses
$b_{kv}$ bytes. Multiplying by $L$ layers gives bytes for one
evaluated token in one request.

Only then multiply by the number of cached tokens across requests.
If lengths differ, use $\sum_r S_r$, not the maximum length times
request count unless you explicitly mean a reservation. The logical
payload and the allocated capacity can be different even before
workspace and allocator metadata are included.

\EngineFigure{foundation-cache-accounting}{Three illustrative requests
each need five token slots. With blocks of four, each reserves eight
slots: five used and three unused. The same model state costs 960
logical bytes or 1536 reserved bytes when each token uses 64 bytes.}
{fig:foundation-cache-accounting}

The figure's unused slots are internal fragmentation. They do not
contain additional context and do not improve the answer. Smaller
blocks can reduce this waste while increasing block-table entries
and management work. The paged-cache chapter will measure that trade;
here the diagram makes the two byte counts visible.

Logical position and physical placement are another pair of distinct
coordinates. Position five is the sixth token. With blocks of four,
it lies in logical block one at offset one. A block table translates
that logical block to whichever physical block the allocator supplied.
Never use physical block order as a position for RoPE.

\section{Turn a traffic count into a time bound}
Suppose an operation must read two billion bytes from a named memory
interface. At an assumed ceiling of one hundred billion bytes per
second, those reads alone need at least 0.02 seconds. Units cancel:
bytes divided by bytes per second gives seconds.

If the same operation needs eight billion FLOPs on an arithmetic path
capable of one trillion FLOPs per second, arithmetic alone needs at
least 0.008 seconds. With ideal overlap, the larger bound dominates.
Adding the two gives a serial model; taking the maximum gives an
ideal-overlap lower bound. They describe different scheduling assumptions.

Neither bound includes every cost automatically. Extra reads, writes,
launches, dependencies and queueing must be counted or acknowledged.
Do not replace a measured time by the bound and call the result a
benchmark. A model is useful because it makes a prediction that an
experiment can challenge.

Arithmetic intensity divides work by traffic at the same boundary.
Eight billion FLOPs divided by two billion bytes gives four FLOP
per byte. If a tile reuses values from local storage, traffic at the
external-memory boundary can fall even though the arithmetic count
stays fixed. Intensity is therefore not a property of the equation
alone; it also depends on how and where it is evaluated.

\section{A rate must name its population}
A batch may produce four output tokens in a 20 ms step. Its aggregate
rate is 200 tokens per second. If those tokens belong to four
independent requests, each request advances by one token in that
step, or 50 tokens per second while the pattern persists.
The aggregate rate is not what one user sees.

Likewise, time per output token is not the same as time to first
token. Prompt processing, queueing and setup contribute differently
to those quantities. Write the start and end events before computing
a rate or an average. The next chapter supplies an explicit trace
and asks which population each statistic summarizes.

Run the arithmetic checks:
\begin{verbatim}
python3 code/foundations/lesson.py units
\end{verbatim}
The output includes 64 bytes per cached token, 960 logical bytes,
1536 reserved bytes and a 20 ms lower bound. These are derived from
illustrative inputs. They are deliberately small enough to recompute
without trusting the program.

\section{Use notation as a debugging tool}
Before accepting an equation, ask four questions. What objects are
being counted? Which axes remain after each operation? What units
cancel? Which assumptions make an equality, approximation or bound
valid? A failed answer usually identifies a missing definition or
a wrong model before it becomes a wrong implementation.

The symbol table is a reference, not a vocabulary test. Return to
it when a later chapter combines weights, KV state and concurrent
requests. Now apply it to an entire small projection, from the shape
of its operands to the assumptions behind its performance model.

\section{Follow one projection through four different counts}
Choose illustrative activations
$\mathbf X\in\mathbb R^{3\times16}$ and weights
$\mathbf W\in\mathbb R^{16\times8}$. Their product
$\mathbf Y=\mathbf X\mathbf W$ has three rows and eight columns.
There are 48 input elements, 128 weights and 24 output elements.
These are three different objects. The input and output sizes grow
with the number of work rows; the weight size does not.

Assume the activations and outputs are stored with two bytes per element.
Now choose a teaching weight format, not a named production format.
Each complete group of sixteen weights stores eight code bytes and
one two-byte scale. Eight groups represent the 128 weights:
\[
 D_W=8(8+2)=80\ \text{bytes},\qquad
 b_w=\frac{80}{128}=\frac58\ \text{bytes/weight}.
\]
The code payload is four bits per weight. Including scales, the effective
storage is five bits per weight. The distinction is not a rounding error:
forgetting the scales would predict 64 bytes instead of 80.
Chapter~\ref{app:formats} followed an actual Q4\_0 format; this different,
explicitly chosen group makes the accounting method visible.

At a hypothetical boundary where the weights and inputs are fetched once
and each output is written once, the payload traffic is
\[
 D=80+(3\cdot16\cdot2)+(3\cdot8\cdot2)
   =224\ \text{bytes}.
\]
With the two-FLOPs-per-multiply--add convention, work is
$F=2\cdot3\cdot16\cdot8=768$ FLOPs. Thus
$I=768/224=24/7\approx3.43$ FLOP/byte.
This counts the projection, not its bias, normalization, scale conversion
or a whole decoder layer. A direct dot product with sixteen multiplications
and fifteen additions has a different literal operation count;
the conventional count is useful only when the reported compute rate
uses the same convention.

\EngineFigure{foundation-projection-accounting}{A derived accounting
route for the illustrative $3\times16$ projection. Shapes determine
elements; the chosen storage format determines bytes. The traffic box
assumes one fetch or store per object at a named boundary. It is not
an observed memory-controller trace.}{fig:foundation-projection-accounting}

The bytes occupied by weights are not automatically the bytes transferred
in an invocation. Reusing a resident weight tile can remove a fetch at
one boundary, while a strided access can transfer extra sectors at another.
Keep a storage column and a traffic column in your notes. When they happen
to be equal, state the reuse assumption that makes them equal.

\section{Make dimensions executable}
An equation can have reasonable-looking numbers and still be meaningless.
Dividing 8 billion FLOPs by 100 billion bytes per second gives
0.08 \emph{FLOP-seconds per byte}, not 0.08 seconds. The divisor describes
the wrong resource. Adding that result to a time does not repair it.

The new lesson stores a quantity as an exact rational value plus a dimension
vector. Its three coordinates count powers of bytes, FLOPs and seconds.
Bytes have vector $(1,0,0)$; bandwidth has $(1,0,-1)$.
Division subtracts those vectors, yielding $(0,0,1)$ for a duration.
This is a small type discipline, not a simulation of a device:
\begin{verbatim}
traffic = Quantity.of(2_000_000_000, "byte")
bandwidth = Quantity.of(100_000_000_000, "byte/s")
seconds = (traffic / bandwidth).in_unit("s")
\end{verbatim}
The value is exactly $1/50$ second. Replacing traffic by FLOPs makes
the final unit check reject the calculation.

Here is the actual division and unit-check path from the lesson:
\begin{minipage}{\linewidth}
\CodeLines{../code/foundations/performance_contract.py}{28}{35}
\end{minipage}
The \texttt{zip} pairs corresponding dimension exponents.
The subtraction changes units; the rational division changes the value.
The final method checks the dimension vector instead of trusting a
variable named \texttt{seconds}. This implementation is deliberately
small; it does not replace a general units library or encode memory
boundaries in its type.

Dimensions alone cannot detect every mistake. The matrix work count and
an unrelated kernel's work count are both FLOPs. A memory controller's
bandwidth and a network link's bandwidth are both bytes per second.
Attaching a boundary, an operation and a precision remains the reader's
job. Unit checking catches a class of errors; semantic checking catches
the wrong objects within the right class.

Decimal and binary prefixes also need explicit conversion. One GiB is
1,073,741,824 bytes. Dividing it by one billion bytes per second gives
1.073741824 seconds, not one second. A gigabit-per-second specification
uses bits, so its nominal byte rate is one eighth as large before
protocol overhead. Name the quantity before using its printed number.
Never compare a traffic estimate in GiB with a bandwidth in GB/s by
silently cancelling the prefixes.

\section{Distinguish payload, reservation and peak residency}
Return to the 64-byte cache footprint per evaluated token. Three requests
with lengths $[1,5,9]$ have fifteen tokens, hence 960 logical bytes.
Private four-token blocks reserve $[4,8,12]$ slots, or 1536 bytes.
Six block-table entries at an illustrative four bytes each add 24 bytes.
The cache and its table therefore occupy 1560 bytes under this narrow
allocation model. Workspace, alignment and allocator objects are still
outside that total.

Multiplying the maximum length by three gives 27 slots, or 1728 bytes.
That is a rectangular maximum-length reservation, not the logical
state of these requests. Neither number is universally wrong.
They describe different allocation policies. The program accepts
individual lengths so that changing one request does not accidentally
change the state attributed to every other request.

Sharing changes the count again. If two requests share an immutable
complete four-token prefix block, its 256 payload bytes need exist only
once, though both requests need a mapping to it. Without sharing they
occupy 512 bytes. An independently writable tail cannot be shared by
the same argument: ownership and modification rules decide whether a
physical object may be counted once. Identical token counts alone do not
establish shareable state.

The 80-byte weight tensor also illustrates alignment. If only its outer
allocation is rounded to a 32-byte multiple, reservation becomes 96 bytes.
If every ten-byte group were independently rounded to 32 bytes, it would
become 256 bytes instead. Those are different physical formats; padding
must follow the actual layout, not a convenient percentage.
If the loader retains the packed 80 bytes while materializing all 128
weights as two-byte values, both buffers coexist for 336 payload bytes.
Their lifetimes determine peak residency. A final checkpoint size says
nothing about whether such temporary duplication happens.

\section{Choose a bound that matches the dependency graph}
The earlier example has a 20 ms memory-resource time and an 8 ms
compute-resource time. If both costs are unavoidable at their named
boundaries and their rates are true ceilings, every implementation
needs at least 20 ms. The maximum is a lower bound even when overlap
is imperfect. It is not a promise that overlap can be implemented.

Suppose a particular two-stage implementation completes a bulk transfer
before allowing its compute stage to begin. Its dependency gives the
stronger serial resource bound $20+8=28$ ms. If independent tiles can
flow through both resources concurrently, a different graph may approach
the maximum after filling its pipeline. First-result latency, steady-state
throughput and pipeline drain are different questions.

\EngineFigure{foundation-bound-dependencies}{Illustrative resource
durations of 20 and 8 ms. The upper schedule permits ideal overlap;
the lower imposes a completion dependency. The resulting 20 and 28 ms
values are conditional models, not timings from a GPU or server.}
{fig:foundation-bound-dependencies}

There is an important distinction between a hardware upper limit and a
measured rate used as a proxy. If a calibration achieved 100 GB/s, another
access pattern might achieve more or less. Dividing by that measured rate
gives an estimate under a transferability assumption, not a rigorous
hardware lower bound. Similarly, a dense matrix peak cannot automatically
price a reduction, a dequantizer or sparse routing. Match the operation
and precision before assigning $\pi$.

To challenge the model, record an observed time and compare it with the
conditional estimate. If a result beats a true bound, inspect your assumptions:
perhaps traffic was served by a cache, work was skipped, the rate was not
a ceiling or the timer stopped at submission rather than completion.
If it is slower, the ratio does not identify the missing cost by itself.
Look for dependencies, small-work inefficiency, additional transfers and
controller work before assigning all the difference to memory bandwidth.

\section{Predict where an improvement stops helping}
With $D/\beta=20$ ms and $F/\pi=8$ ms, halving traffic changes the
ideal-overlap model from 20 to 10 ms. Halving it again gives a 5 ms
memory term, but the result is now bounded by the unchanged 8 ms compute
term. Four times fewer bytes did not produce four times less step time.
The limiting resource moved. This is the reason to keep both terms
visible rather than extrapolating a bandwidth improvement indefinitely.

Nor does an isolated step determine end-to-end benefit. Choose an
illustrative request with 20 ms in an optimizable stage and 20 ms elsewhere.
Halving the stage gives a 30 ms request instead of 40 ms:
the speed ratio is $40/30=4/3$, not two. For an original optimizable
fraction $f$ and its speed ratio $r$, a serial, otherwise unchanged model
gives
\[
 \text{speed ratio}=\frac{1}{(1-f)+f/r}.
\]
Here $f=1/2$ and $r=2$. This algebra describes the stated decomposition;
it does not assume independent speed ratios can be multiplied or that
queueing stays unchanged when a service becomes faster.

A batch adds another population. Under the earlier 20 ms step bound,
four output tokens imply an aggregate ceiling of 200 tokens/s.
For four one-token-per-step requests, each advances at at most 50 tokens/s
under the same model. Speculation may emit several accepted tokens for
one request; padding may emit no useful token for a slot. Count accepted,
useful outputs instead of equating launched work rows with delivered tokens.

\section{Run, explain and deliberately break the accounting}
Run the new lesson without downloading a model:
\begin{verbatim}
python3 code/foundations/performance_contract.py
python3 -m unittest discover -s code/foundations
\end{verbatim}
The output reports 224 projection bytes, 768 FLOPs, cache payload and
reservation separately, and exact 20 ms and 28 ms conditional times.
The tests enumerate tiny projection loops independently of the closed
work formula. Cache tests append padding slots explicitly before comparing
with ceiling division. These alternative calculations make off-by-one
and omitted-factor mistakes visible.

In a scratch copy, omit scale metadata or multiply weight bytes by the
work-row count. The first mutation changes the 80-byte fixture to 64;
the second incorrectly duplicates the model. Swapping bandwidth and compute
ceilings is rejected by dimensions. None of these tests establishes that
the chosen memory boundary is correct for an industrial kernel.
That requires implementation inspection and measurement.

\begin{ConstraintLedger}
Memory bandwidth & moves & Name the boundary and reuse before converting traffic to time.\\
Memory capacity & spends & Metadata, padding and overlapping buffer lifetimes add to payload.\\
Compute & moves & FLOP counts need an operation-specific rate and counting convention.\\
Latency & risks & Overlap, dependencies and queueing distinguish a bound from a prediction.\\
Correctness & buys & Unit and independent counting checks reject plausible accounting errors.\\
\end{ConstraintLedger}

You can now derive a prediction and identify the conditions that might
falsify it. The next chapter turns that prediction into an experiment
whose output, timestamps and failures remain inspectable.
\section{Worked problems}
\subsection*{1. A packed batch}
\textbf{Problem.} Five rows of width four multiply a $4\times3$ matrix.
Give the output shape and weight count.
\textbf{Solution.} Output shape is $5\times3$ and there are twelve
matrix weights. The batch does not duplicate model parameters.
\subsection*{2. Wasted reservation}
\textbf{Problem.} Reproduce the unused KV capacity for three five-token
requests, four-token blocks and 64 bytes per token.
\textbf{Solution.} Each reserves eight slots and wastes three.
Total unused capacity is $3\cdot3\cdot64=576$ bytes.
\subsection*{3. Two rates}
\textbf{Problem.} Four requests each emit one token every 20 ms.
What are the aggregate and per-request rates?
\textbf{Solution.} Aggregate rate is $4/0.02=200$ tokens/s.
Each request sees $1/0.02=50$ tokens/s, excluding other delays.
\subsection*{4. Count the format, not only the codes}
\textbf{Problem.} A $16\times8$ matrix uses the teaching format of
sixteen weights per group, eight code bytes and a two-byte scale.
Find payload bytes, effective bits per weight and traffic for three
input rows with two-byte inputs and outputs.
\textbf{Solution.} Eight groups cost $8(8+2)=80$ bytes.
Effective precision is $8\cdot80/128=5$ bits/weight.
Traffic under one-fetch/store assumptions is $80+96+48=224$ bytes.
\Needspace{10\baselineskip}
\subsection*{5. Two valid cache totals}
\textbf{Problem.} Requests have lengths $[1,5,9]$, use 64 bytes per token
and private four-token blocks. Table entries use four bytes. Find
logical bytes and reservation including the table.
\textbf{Solution.} Logical bytes are $(1+5+9)64=960$.
There are $1+2+3=6$ blocks, reserving $6\cdot4\cdot64=1536$ bytes.
The table adds 24 bytes, giving 1560. This excludes other allocator costs.
\subsection*{6. The new limiting resource}
\textbf{Problem.} Resource times are 20 ms for memory and 8 ms for
compute. Reduce traffic by four without changing work or rates.
What is the ideal-overlap bound? What if an unrelated 20 ms serial
stage originally followed a 20 ms stage that is sped up by two?
\textbf{Solution.} Memory falls to 5 ms, so the bound is 8 ms,
not 5. The separate end-to-end example changes from 40 to 30 ms,
a $4/3$ speed ratio. These are conditional derived models.

